\documentclass[10pt,a4paper]{amsart}
\usepackage{amsmath,amssymb,amsthm,graphicx}
% For commutative diagrams we use Xy-pic
\usepackage[all]{xy}

% We use 2cell for 2-commutative diagrams.
\xyoption{2cell}
% Theorem environments.
%
\theoremstyle{plain}
\newtheorem{theorem}[subsection]{Theorem}
\newtheorem{proposition}[subsection]{Proposition}
\newtheorem{lemma}[subsection]{Lemma}

\theoremstyle{definition}
\newtheorem{definition}[subsection]{Definition}
\newtheorem{example}[subsection]{Example}
\newtheorem{exercise}[subsection]{Exercise}
\newtheorem{situation}[subsection]{Situation}

\theoremstyle{remark}
\newtheorem{remark}[subsection]{Remark}
\newtheorem{remarks}[subsection]{Remarks}

\numberwithin{equation}{subsection}

% Macros
%
\def\lim{\mathop{\mathrm{lim}}\nolimits}
\def\colim{\mathop{\mathrm{colim}}\nolimits}
\def\Spec{\mathop{\mathrm{Spec}}}
\def\Hom{\mathop{\mathrm{Hom}}\nolimits}
\def\Ext{\mathop{\mathrm{Ext}}\nolimits}
\def\SheafHom{\mathop{\mathcal{H}\!\mathit{om}}\nolimits}
\def\SheafExt{\mathop{\mathcal{E}\!\mathit{xt}}\nolimits}
\def\Sch{\mathit{Sch}}
\def\Mor{\mathop{\mathrm{Mor}}\nolimits}
\def\Ob{\mathop{\mathrm{Ob}}\nolimits}
\def\Sh{\mathop{\mathit{Sh}}\nolimits}
\def\NL{\mathop{N\!L}\nolimits}
\def\CH{\mathop{\mathrm{CH}}\nolimits}
\def\proetale{{pro\text{-}\acute{e}tale}}
\def\etale{{\acute{e}tale}}
\def\QCoh{\mathit{QCoh}}
\def\Ker{\mathop{\mathrm{Ker}}}
\def\Im{\mathop{\mathrm{Im}}}
\def\Coker{\mathop{\mathrm{Coker}}}
\def\Coim{\mathop{\mathrm{Coim}}}

% Boxtimes
%
\DeclareMathSymbol{\boxtimes}{\mathbin}{AMSa}{"02}

%
% Macros for moduli stacks/spaces
%
\def\QCohstack{\mathcal{QC}\!\mathit{oh}}
\def\Cohstack{\mathcal{C}\!\mathit{oh}}
\def\Spacesstack{\mathcal{S}\!\mathit{paces}}
\def\Quotfunctor{\mathrm{Quot}}
\def\Hilbfunctor{\mathrm{Hilb}}
\def\Curvesstack{\mathcal{C}\!\mathit{urves}}
\def\Polarizedstack{\mathcal{P}\!\mathit{olarized}}
\def\Complexesstack{\mathcal{C}\!\mathit{omplexes}}
% \Pic is the operator that assigns to X its picard group, usage \Pic(X)
% \Picardstack_{X/B} denotes the Picard stack of X over B
% \Picardfunctor_{X/B} denotes the Picard functor of X over B
\def\Pic{\mathop{\mathrm{Pic}}\nolimits}
\def\Picardstack{\mathcal{P}\!\mathit{ic}}
\def\Picardfunctor{\mathrm{Pic}}
\def\Deformationcategory{\mathcal{D}\!\mathit{ef}}
\usepackage{cleveref}
\crefname{originallemma}{original lemma}{original lemmas}
\Crefname{originallemma}{Original lemma}{Original lemmas}
\makeatletter
\newlabel{sites-lemma-back-and-forth}{{00XT}{source}{}{sourcepointer-1}{}}
\newlabel{sites-lemma-back-and-forth@cref}{{[originallemma][0][]00XT}{[1][0][]source}{}{}{}}
\makeatother
\title{For a monomorphism the unit of the big pullback-pushforward adjunction is an isomorphism}
\author{The Stacks Project Authors}
\date{stacks-project:etale-cohomology.tex:local-snapshot-20261011; GNU FDL 1.2 or later}
\begin{document}
\maketitle
\noindent\textit{Editorial note.} All delivered mathematical text is the author's own, copied from the pinned Stacks Project chapter file etale-cohomology.tex; no statement or proof step was rewritten, reordered or removed.\par\medskip
\noindent\textit{Editorial note.} The proof is the author's complete argument: it lists the properties of the induced functor and names the criterion that turns them into the conclusion.\par\medskip
\noindent\textit{Editorial note.} The label lemma-monomorphism-big-push-pull occurs in no other staged Stacks chapter, so no cross-chapter restatement had to be weighed for this card.\par\medskip
\noindent\textit{Editorial note.} The single labelled reference occurring inside the counted spans points to a chapter of the Project not staged in this card's source tree; it is carried in assembly.external\_references as a non-load-bearing navigation pointer recorded at the invoking line.\par\medskip
\noindent\textit{Editorial note.} The preamble is the chapter's own native preamble excerpt already used for this source: the Xy-pic setup (native preamble lines 14-18) and the native theorem declarations with the mathematical macros (native preamble lines 154-217).\par\medskip
\noindent\textit{Allowed standard prerequisites (editorial).}\begin{itemize}
\item Big site (Sch/X)\_tau and the morphism of sites induced by a morphism of schemes. Same chapter, the big Zariski, etale, smooth, syntomic and fppf sites of a scheme and the functor (Sch/X)\_tau -> (Sch/Y)\_tau induced by f. Supplies the functor whose properties are checked and the morphism of topoi f\_big whose adjunction unit is examined.
\item Monomorphism of schemes and the induced functor on big sites. Same chapter, the fact that for a monomorphism f the functor on the big sites is fully faithful. Supplies the full faithfulness which, with continuity and cocontinuity, is the hypothesis of the cited criterion.
\item Criterion for the adjunction unit to be an isomorphism. Sites, the result that for a fully faithful continuous and cocontinuous functor the canonical map f\textasciicircum{}\{-1\}f\_*F -> F is an isomorphism. Supplies the conclusion applied to the big morphism f\_big.
\end{itemize}\medskip
\ifdefined\hypertarget\hypertarget{sourcepointer-1}{}\fi
\noindent\textit{Source reference.} Original lemma 00XT: sources/stacks-etale-cohomology/pinned-original-source.tex, source lines 6837--6837. Cross-chapter navigation pointer invoked at proof line 6837, where the author uses the criterion that a fully faithful continuous and cocontinuous functor gives an isomorphism for the unit of the adjunction. The target chapter is not staged in this card's source tree, so the pointer records the invoking line; it only resolves the reference for the reader, and no bibliography citation is involved.\par\medskip
\begin{lemma}
\label{lemma-monomorphism-big-push-pull}
Let $\tau \in \{Zariski, \etale, smooth, syntomic, fppf\}$.
Let $f : X \to Y$ be a monomorphism of schemes.
Then the canonical map
$f_{big}^{-1}f_{big, *}\mathcal{F} \to \mathcal{F}$
is an isomorphism for any sheaf $\mathcal{F}$ on
$(\Sch/X)_\tau$.
\end{lemma}

\begin{proof}
In this case the functor $(\Sch/X)_\tau \to (\Sch/Y)_\tau$
is continuous, cocontinuous and fully faithful. Hence the result follows from
Sites, Lemma \ref{sites-lemma-back-and-forth}.
\end{proof}

\end{document}
